\section{The complete auxiliary tuple}\label{app:witnesses}

This appendix fixes one finite expansion of the systems used in the proof.
Its purpose is to make explicit which coordinates the single-fold assertion
counts. Each call has a fresh internal tuple. Its displayed output is
counted by its caller, and its displayed inputs are already available.
Equal-valued coordinates in distinct calls remain distinct variables.
Every coordinate below ranges over \(\NN\).

\paragraph{Primitive conventions.}
A polynomial equality introduces no auxiliary variable. A weak comparison
\(F\ge G\) introduces the single slack \(F-G\); a strict comparison
\(F>G\) introduces \(F-G-1\). Naming a polynomial expression introduces
one variable with its defining equality. Rational comparisons first clear
their positive denominators, so their slacks are integer differences.
For a disjoint case construction,
include every Boolean selector, every local guard slack, and every
coordinate of every local call. Multiply all local equations by the case
selector and set all local coordinates to zero when that case is inactive.
These rules apply recursively to nested calls. In particular, a zeroed
inactive rational code has no ungated positive-denominator requirement.

\paragraph{Rational and point arithmetic.}
The six coordinates of a rational code are
\((p,n,d,t,s,q)\) from \eqref{eq:rational-code}: the two sign parts,
positive denominator, inverse residue, its interval slack, and the inverse
quotient. A point has its identity flag and both complete rational codes.
For an addition, the caller counts the output point once. The internal
addition tuple consists of seven selectors, four local
six-coordinate slope codes (Cases 3, 4, 6, and 7), and the four strict
comparison slacks in those cases. Other guards are equalities and require
no further witnesses. Clearing denominators introduces no free multiplier.
A maximum uses two selectors and one local comparison slack in each case;
the first case receives a tie. All inactive local coordinates are zero.

Fix one addition chain for \([m_*]Q\) once and for all; for example, use
\(Q,2Q,\ldots,m_*Q\). Each addition contributes its output point and the
internal addition tuple just described. A height call \(W(Q)\) has one
displayed height output, counted by its caller. Its internal tuple contains
the scalar chain and the maximum tuple applied to the final x-coordinate's
absolute numerator and denominator. The chain length is a fixed constant.

\paragraph{The growth graph.}
For a representative test on a point \(Q\), retain its height output
\(h_Q\), the internal height tuple, the four codes of \(Q+t\) and their
addition tuples, and six comparison slacks: one for \(y(Q)\ge0\), one for
\(h_Q>1\), and four for \(x(Q)\ge x(Q+t)\). Affineness is a flag equation.
This test has unique internal witnesses for a fixed representative; the
bracket and window are what fix the representative itself.

Write \(\xi\) for the entire internal tuple of a call \(v=g(u)\).
It consists of the following blocks:
\begin{itemize}
\item The codes of \(Q,R\) and both representative-test tuples, including
      their height outputs \(h_Q,h_R\).
\item The codes of \(Q+P,Q-P,R+P,R-P,[2]Q\), all five addition tuples,
      and their five height outputs with five complete internal height tuples.
\item The two neighbor-height maxima with both complete maximum tuples,
      and \(L_0,D_0\) with their product equalities.
\item The four slacks in the bracket and window: if the neighbor maxima
      are \(N_Q,N_R\), these are
      \[
       u+W_1-h_Q,\quad N_Q-u-W_1-1,\quad
       \kappa D_0-L_0,\quad \kappa L_0-D_0.
      \]
\end{itemize}
Impose \(v=h_R\). Here the already present internal coordinate \(h_R\)
and the displayed output \(v\) are distinct variables fixed to the same
value. The heights \(h_Q,h_R\) are reused in the products and inequalities.
All signs, denominators, slopes, selectors, and scalar-chain points remain
inside \(\xi\). The integer indices \(k,r\) and the real quantities in
the growth proof are not variables of this tuple.

\paragraph{Powers and binary containment.}
For displayed arguments \((b,n,c)\), one complete internal power tuple is
\[
 (u,v,X,Y,F,H,q,q',r,\sigma_b,\sigma_F,\sigma_{Fv},
       \sigma_H,\sigma_{Hv},\sigma_r,\xi),
\]
where the six slacks, in order, have values
\[
 b-2,\quad F-1,\quad v-F-1,\quad
 H-1,\quad v-H-1,\quad H-r-1.
\]
The output \(c\) is counted by the caller. No sign variable accompanies
the natural square roots \(X,Y\); no Pell index is quantified.

For a binary-containment call on \((x,y)\), name \(e=y+1\) and
\(A=B+1\). Include
\[
 (e,A,B,H,N,s,d,r,t,\sigma_{xy},\sigma_d,\sigma_r)
\]
and three separate complete internal power tuples for
\(\Pow(2,e,B)\), \(\Pow(B,x,H)\), and \(\Pow(A,y,N)\).
The three slacks are \(y-x,B-d-1,H-r-1\). The quotients and remainders
are precisely \(s,d,r\), and the parity quotient is \(t\).
These calls, including all of their growth tuples, remain present when
\(x=0\) or \(y=0\).

\paragraph{The computation tuple.}
Let the fixed machine have \(s\) states, \(r\) counters, and
\(\tau\) test states. Enumerate each of these finite sets once.
The displayed machine coordinates are
\[
 (w,B,b,k,H,J,M),\qquad
 (I_q)_q,\qquad (Z_q,N_q,V_q)_{q\text{ a test state}},
 \qquad (R_j,\rho_j,\eta_j)_{0\le j\le r}.
\]
Include the \(r+5\) comparison slacks
\[
 B-2C-1,\quad k-1,\quad (J-\eta_j-1)_{0\le j\le r},
 \quad \rho_0-b,\quad B-2\rho_0-1.
\]
Name \(D=(2b-1)M\) and, for each test state, \(D_q=(B-1)V_q\),
using their defining equalities. Include the two internal power tuples
for \(\Pow(2,w,B)\) and \(\Pow(B,k,H)\). Finally include a complete
binary-containment tuple for each of
\[
 \begin{gathered}
  I_q\bits M,\qquad Z_q\bits I_q,\qquad N_q\bits I_q,\\
  R_j\bits D\quad(0\le j\le r),\qquad
  R_{j(q)}\bits D_q\quad(q\text{ a test state}),
 \end{gathered}
\]
where \(j(q)\) is the counter tested at \(q\). There are
\(s+3\tau+r+1\) such calls. Arc masks are expressions in the displayed
state and test masks and introduce no variables. The expression
\(C=1+s+\sum_jd_j\) also introduces no variable.

This is a finite recursive specification of the whole auxiliary tuple.
Its length depends only on the fixed machine and fixed arithmetic
constructions, not on an input or its halting time. Proof digits,
configurations, and time indices are not extra coordinates.
The computation proof fixes the actual run first, then its base,
then every displayed value above. The preceding graphs fix all their
internal coordinates, including every inactive branch. Replacing the
resulting finite list of equations by its sum of squares introduces no
coordinate and preserves exactly this complete tuple.
